% !TEX root = ../elteikthesis_en.tex

\chapter{User Documentation}
\label{ch:user}

This chapter describes how to install, configure, and use the ELTE Informatics Chatbot.
The system consists of a one-time data pipeline that builds the knowledge base, and a
FastAPI server that answers questions through a browser-based interface or a command-line
client.


% ─────────────────────────────────────────────────────────────────────────────
\section{System Requirements}
\label{sec:user-requirements}

The application is designed to run on a single personal computer without a GPU.
Table~\ref{tab:sys-req} lists the minimum requirements.

\begin{table}[H]
    \centering
    \begin{tabular}{ | m{0.28\textwidth} | m{0.62\textwidth} | }
        \hline
        \textbf{Component} & \textbf{Requirement} \\
        \hline \hline
        Operating system & Windows 10/11, macOS 12+, or Ubuntu 22.04+ \\
        \hline
        RAM & 8~GB (16~GB recommended) \\
        \hline
        Disk space & $\approx$6~GB (models + vector store + crawled data) \\
        \hline
        Python & 3.10 or newer \\
        \hline
        Ollama & Latest release (\url{https://ollama.com}) \\
        \hline
        Internet access & Required only during the initial setup phase \\
        \hline
    \end{tabular}
    \caption{Minimum system requirements}
    \label{tab:sys-req}
\end{table}


% ─────────────────────────────────────────────────────────────────────────────
\section{Installation}
\label{sec:user-install}

\subsection{Install Ollama}

Download and install Ollama from \url{https://ollama.com/download} for your operating
system. After installation, start the Ollama background service:

\begin{lstlisting}[language=bash]
ollama serve
\end{lstlisting}

Pull the two language models used in the system:

\begin{lstlisting}[language=bash]
ollama pull llama3.2:3b
ollama pull gemma3:4b
\end{lstlisting}

\subsection{Clone the Repository and Create a Virtual Environment}

\begin{lstlisting}[language=bash]
git clone https://github.com/borohull/Thesis.git
cd Thesis/elte_chat
python -m venv .venv
\end{lstlisting}

Activate the virtual environment:

\begin{compactitem}
    \item \textbf{Windows:} \texttt{.venv\textbackslash Scripts\textbackslash activate}
    \item \textbf{macOS / Linux:} \texttt{source .venv/bin/activate}
\end{compactitem}

\bigskip
Install the Python dependencies:

\begin{lstlisting}[language=bash]
pip install -r requirements.txt
\end{lstlisting}


% ─────────────────────────────────────────────────────────────────────────────
\section{Building the Knowledge Base}
\label{sec:user-pipeline}

Before the chatbot can answer questions, the knowledge base must be built.
This is a one-time process consisting of three steps.

\subsection{Step 1 — Crawl the ELTE Websites}

The crawler downloads HTML pages and PDF documents from the ELTE Faculty of Informatics
website and saves them to \texttt{data/raw/}:

\begin{lstlisting}[language=bash]
python scripts/crawler.py
\end{lstlisting}

To start fresh and discard previously downloaded files, pass the \texttt{-{}-reset} flag:

\begin{lstlisting}[language=bash]
python scripts/crawler.py --reset
\end{lstlisting}

The crawler performs a breadth-first traversal starting from the faculty homepage.
Downloaded files are tracked in a manifest so that subsequent runs only fetch new or
changed documents.

\subsection{Step 2 — Chunk and Embed Documents}

Open JupyterLab or Jupyter Notebook and run the following notebooks in order:

\begin{enumerate}
    \item \textbf{\texttt{notebooks/01\_data\_ingestion.ipynb}} — loads the raw files,
          splits them into overlapping text chunks, and writes the result to
          \texttt{data/processed/chunks.json}.
    \item \textbf{\texttt{notebooks/02\_embedding.ipynb}} — encodes each chunk with the
          \texttt{all-MiniLM-L6-v2} sentence-transformer model and stores the resulting
          vectors in the ChromaDB collection \texttt{elte\_ik} under
          \texttt{data/processed/chroma\_db/}.
\end{enumerate}

Both notebooks must be run to completion before starting the API server.
Running notebook~02 on a CPU with several thousand chunks typically takes a few minutes.

\subsection{Configuration}

All tunable parameters are defined in \texttt{app/settings.py}.
The most commonly adjusted settings are listed in Table~\ref{tab:settings}.

\begin{table}[H]
    \centering
    \begin{tabular}{ | m{0.3\textwidth} | m{0.18\textwidth} | m{0.42\textwidth} | }
        \hline
        \textbf{Setting} & \textbf{Default} & \textbf{Description} \\
        \hline \hline
        \texttt{OLLAMA\_MODEL} & \texttt{llama3.2:3b} & Default language model \\
        \hline
        \texttt{TOP\_K} & \texttt{5} & Chunks retrieved per query \\
        \hline
        \texttt{TEMPERATURE} & \texttt{0.1} & LLM sampling temperature \\
        \hline
        \texttt{EMBEDDING\_MODEL} & \texttt{all-MiniLM-L6-v2} & Must match notebook~02 \\
        \hline
    \end{tabular}
    \caption{Key configuration parameters in \texttt{app/settings.py}}
    \label{tab:settings}
\end{table}

\noindent\textbf{Important:} \texttt{EMBEDDING\_MODEL} must not be changed after
notebook~02 has been run. The query encoder and the document index must use the same
model; a mismatch produces incorrect retrieval results.


% ─────────────────────────────────────────────────────────────────────────────
\section{Starting the Server}
\label{sec:user-server}

Run the following command from the project root (with the virtual environment active and
Ollama running):

\begin{lstlisting}[language=bash]
uvicorn app.main:app --reload
\end{lstlisting}

The server starts on \texttt{http://localhost:8000}.
The \texttt{--reload} flag restarts the server automatically when source files change;
omit it in production.

To verify the server is healthy and Ollama is reachable, open a browser or run:

\begin{lstlisting}[language=bash]
curl http://localhost:8000/health
\end{lstlisting}

A successful response looks like:

\begin{lstlisting}
{"status": "ok", "ollama": true}
\end{lstlisting}

If \texttt{"ollama": false} appears, Ollama is not running or the selected model has
not been pulled.


% ─────────────────────────────────────────────────────────────────────────────
\section{Using the Web Interface}
\label{sec:user-web}

Open \texttt{http://localhost:8000/static/index.html} in any modern browser.
The interface consists of three areas: a left sidebar for chat history, a central chat
panel, and a right sidebar for model selection.

\subsection{Asking a Question}

Type a question in the input field at the bottom of the chat panel and press
\textbf{Enter} or click the send button. The assistant retrieves the most relevant
document chunks from the knowledge base, constructs a prompt, and returns an answer.

Questions should relate to the ELTE Faculty of Informatics — for example:

\begin{compactitem}
    \item \emph{What is a strong prerequisite?}
    \item \emph{Who is the Dean of the Faculty of Informatics?}
    \item \emph{Can I take a follow-up subject before completing its prerequisite?}
\end{compactitem}

For questions outside the scope of the knowledge base, the system responds that the
information is not available in its context rather than fabricating an answer.

\subsection{Selecting a Language Model}

The right sidebar lists all Ollama models currently available on the machine.
Click a model name to switch to it for the current session.
The default model is \texttt{llama3.2:3b}; \texttt{gemma3:4b} is also available after
running \texttt{ollama pull gemma3:4b}.

Response latency depends on model size and hardware. On a CPU-only machine with 8~GB
of RAM, typical response times are 30--80~seconds per query.

\subsection{Managing Chat Sessions}

Each conversation is stored as a named session. Sessions appear in the left sidebar and
persist across browser refreshes. To start a new conversation, click \textbf{New chat}.
To delete a session, hover over its entry in the sidebar and click the
\texttimes\ button that appears.

\subsection{Uploading Custom Documents}

In addition to the crawled ELTE content, users can upload their own PDF, HTML, or DOCX
files to extend the knowledge base. Click the \textbf{Upload document} button above the
input field, select a file (maximum 20~MB), and wait for the confirmation message.
Uploaded documents are chunked, embedded, and added to the vector store immediately;
they are available for retrieval in the same session without restarting the server.


% ─────────────────────────────────────────────────────────────────────────────
\section{Using the Command-Line Client}
\label{sec:user-cli}

A lightweight CLI client is provided for scripting and testing without a browser.
The FastAPI server must be running before the client is started:

\begin{lstlisting}[language=bash]
python scripts/chat_cli.py
\end{lstlisting}

Type a question at the \texttt{>} prompt and press Enter. The client sends the question
to the server and prints the answer together with the source references.
Type \texttt{exit} or press \texttt{Ctrl+C} to quit.


% ─────────────────────────────────────────────────────────────────────────────
\section{Interaction Logs}
\label{sec:user-logs}

Every question and answer is automatically recorded in a SQLite database at
\texttt{data/logs/chat\_logs.db}. Each row stores the timestamp, the user's message,
the generated answer, the retrieved source references, the end-to-end response time in
milliseconds, and any error that occurred. Logs can be inspected with any SQLite viewer
or via the optional cell in \texttt{notebooks/04\_evaluation.ipynb}.
